\chapter{Struttura Comune: BaseAgent}

Tutti e tre i tipi di agente ereditano da \texttt{BaseAgent}, che implementa le funzionalit\`a
condivise: percezione dell'ambiente, ricezione e invio messaggi, movimento tramite A*,
gestione dell'energia e rilevamento dello stato di stallo.

\section{Percezione e Mappa Locale}

Ogni agente mantiene una \textit{mappa locale} rappresentata come matrice numpy, inizialmente
tutta \texttt{UNKNOWN}. Ad ogni passo, la fase \texttt{STAGE\_SENSE} aggiorna le celle entro
il raggio di visione con il tipo reale della griglia, applicando la logica di \textit{ray
casting}: le celle dietro un ostacolo non vengono aggiornate anche se geometricamente entro
il raggio, simulando la visibilit\`a lineare degli agenti fisici. Gli oggetti nella zona di
visione vengono aggiunti alla lista \texttt{known\_objects}; le celle magazzino vengono
aggiunte a \texttt{known\_warehouses}.

\section{Integrazione della Mappa Condivisa}

Alla ricezione di un \texttt{MapDataMessage}, l'agente aggiorna le celle della propria
mappa ancora \texttt{UNKNOWN} con i dati del mittente. Il merge estrae anche le celle
magazzino, popolando \texttt{known\_warehouses}: in questo modo i Retriever scoprono
le stazioni di ricarica ben prima di visitarle, purché uno Scout o un Coordinator
vicino le abbia già esplorate.

\section{Protocollo ClearWay}

Quando un agente calcola un percorso A* verso il proprio obiettivo ma la cella successiva
\`e occupata da un altro agente fermo, anzich\'e attendere o ricalcolare un percorso
alternativo, invia un \texttt{ClearWayMessage} direttamente all'occupante. L'agente
ricevente, se libero da vincoli immediati, si sposta sulla cella adiacente pi\`u vicina
libera. Questa negoziazione \`e ricorsiva (\texttt{depth} nel messaggio) fino a una
profondità massima: se la catena di spostamenti non si risolve, l'agente iniziante
rilascia il messaggio e lascia che A* trovi una via alternativa al passo successivo.
Il meccanismo riduce significativamente i blocchi in corridoi stretti senza richiedere
alcuna struttura di prenotazione globale.

\section{Rilevamento dello Stallo}

L'agente registra lo storico delle posizioni recenti. Se per un certo numero di passi
consecutivi (\texttt{stuck\_threshold}) la posizione non cambia, il campo
\texttt{stuck\_counter} viene incrementato. Superata la soglia, l'agente esegue una
\textit{stuck recovery}: invalida il percorso corrente, aggiunge le posizioni degli
agenti vicini come evitate e ricalcola un percorso alternativo verso una destinazione
leggermente ruotata. Se lo stallo persiste, l'agente si sposta temporaneamente verso
una cella casuale libera per sbloccarsi.

Quando il rilevatore di loop o di stallo abbandona un target, la posizione viene inserita
nella blacklist \texttt{unreachable\_targets} con il passo corrente: in questo modo il
ciclo decisionale dell'agente non riassegna immediatamente lo stesso obiettivo irraggiungibile
nel passo successivo.
